%============================================================
% CH1 · LESSON 3 — STORING A GRAPH
%============================================================
\subsection{Storing a Graph}
\label{ssec:ch1-storing}

A drawing is for humans; a program needs the graph in memory, in a form that answers the questions the algorithm asks. Almost every algorithm asks one question over and over: \textbf{``who are the neighbours of $u$?''} The structure that answers it in time proportional to the answer --- the \emph{adjacency list} --- is what you will use in 95\% of problems. This lesson shows it together with the three alternatives, so that you know when each one is the right tool.

Throughout the lesson we store the same weighted undirected graph $G_2$:

\begin{diagrambox}
\begin{tikzpicture}[x=1.4cm, y=1.1cm]
    \node[gv] (1) at (0,0.8) {1};
    \node[gv] (2) at (0,-0.8) {2};
    \node[gv] (3) at (1.4,0) {3};
    \node[gv] (4) at (2.9,0) {4};
    \node[gv] (5) at (4.4,0) {5};
    \draw[ge] (1) -- node[gw, left] {4} (2);
    \draw[ge] (1) -- node[gw, above right] {1} (3);
    \draw[ge] (2) -- node[gw, below right] {2} (3);
    \draw[ge] (3) -- node[gw, above] {7} (4);
    \draw[ge] (4) -- node[gw, above] {3} (5);
    \node[draw=darkgray, fill=lightgray!60, rounded corners=2pt, font=\ttfamily\small, align=left, inner sep=4pt, anchor=west] at (5.4,0)
        {5 5\\1 2 4\\1 3 1\\2 3 2\\3 4 7\\4 5 3};
    \node[font=\scriptsize\color{darkgray}] at (6.0,1.35) {input};
\end{tikzpicture}
\figcap{$G_2$: $n=5$ vertices, $m=5$ weighted edges, given as ``$u$ $v$ $w$'' lines.}
\end{diagrambox}

%------------------------------------------------------------
\subsubsection*{1. Edge list --- just keep the input}

The simplest representation is the input itself: an array of $m$ triples $(u,v,w)$.

\begin{diagrambox}
\begin{tikzpicture}
    \dpindex{ei}{1.4}{$i$}{0}{4}
    \dparray{eu}{0.7}{$u$}{1/cell, 1/cell, 2/cell, 3/cell, 4/cell}
    \dparray{ev}{0}{$v$}{2/cell, 3/cell, 3/cell, 4/cell, 5/cell}
    \dparray{ew}{-0.7}{$w$}{4/compcell, 1/compcell, 2/compcell, 7/compcell, 3/compcell}
    \node[note, anchor=west, text width=5.2cm] at (5.2,0.35) {Good for: algorithms that look at every edge as a whole --- sorting edges by weight (Kruskal, Ch.~10), relaxing all edges (Bellman--Ford, Ch.~8), DSU (Ch.~6), counting degrees (Lesson~\ref{ssec:ch1-vocabulary}).\\[2pt] Bad for: ``neighbours of $u$'' costs $\bigO{m}$.};
\end{tikzpicture}
\figcap{The edge list of $G_2$: three parallel arrays (or one \texttt{vector<array<int,3>>}).}
\end{diagrambox}

\begin{lstlisting}[style=snippet]
struct Edge { int u, v, w; };
vector<Edge> edges(m);
for (auto &e : edges) cin >> e.u >> e.v >> e.w;
sort(edges.begin(), edges.end(), [](const Edge &a, const Edge &b) { return a.w < b.w; });
\end{lstlisting}

%------------------------------------------------------------
\subsubsection*{2. Adjacency matrix --- one cell per pair}

An $n\times n$ table with \texttt{g[u][v]} $=$ weight of the edge $u$--$v$ (or 1/0 for ``edge / no edge'' in an unweighted graph). An undirected graph gives a \emph{symmetric} matrix.

\begin{diagrambox}
\begin{tikzpicture}[x=0.8cm, y=0.65cm]
    \foreach \j in {1,...,5} \node[hdr, minimum width=7mm, minimum height=5.5mm] at (\j,0.9) {\j};
    \foreach \i in {1,...,5} \node[hdr, minimum width=7mm, minimum height=5.5mm] at (0,1-\i-0.1) {\i};
    \foreach \row [count=\i] in {{0,4,1,0,0},{4,0,2,0,0},{1,2,0,7,0},{0,0,7,0,3},{0,0,0,3,0}}{
        \foreach \val [count=\j] in \row {
            \ifnum\val=0
                \node[emptycell, minimum width=7mm, minimum height=5.5mm] at (\j,1-\i-0.1) {$\cdot$};
            \else
                \node[compcell, minimum width=7mm, minimum height=5.5mm, font=\small\bfseries] at (\j,1-\i-0.1) {\val};
            \fi
        }
    }
    \node[font=\scriptsize\color{darkgray}] at (-0.1,0.9) {$u\backslash v$};
    \draw[accentorange, line width=1.2pt, rounded corners=2pt] (0.55,-2.45) rectangle (5.45,-1.75);
    \node[note, anchor=west, text width=5.4cm] at (6.2,-1.6) {Row 3 lists the neighbours of vertex 3 --- but you must scan \emph{all} $n$ cells to find the three non-empty ones.\\[3pt]
        Good for: $n\le$ a few thousand, dense graphs, ``is $u$--$v$ an edge?'' in $\bigO{1}$, Floyd--Warshall (Ch.~9).\\[2pt]
        Bad for: memory $n^2$; neighbours cost $\bigO{n}$; parallel edges collapse into one cell.};
\end{tikzpicture}
\figcap{Adjacency matrix of $G_2$ (dots are ``no edge''). It is symmetric about the diagonal because $G_2$ is undirected.}
\end{diagrambox}

\begin{lstlisting}[style=snippet]
vector<vector<int>> g(n + 1, vector<int>(n + 1, 0));      // n <= ~5000 only!
for (int i = 0; i < m; i++) {
    int u, v, w; cin >> u >> v >> w;
    g[u][v] = g[v][u] = w;          // parallel edges: keep min(w) if needed
}
\end{lstlisting}

%------------------------------------------------------------
\subsubsection*{3. Adjacency list --- the default}

For each vertex keep a dynamic array of its neighbours: \texttt{vector<vector<int>> adj(n + 1)}. An undirected edge $u$--$v$ is pushed twice, as $v$ into \texttt{adj[u]} and as $u$ into \texttt{adj[v]}. For a weighted graph each entry is a pair (neighbour, weight): \texttt{vector<vector<pair<int,int>>>}.

\begin{diagrambox}
\begin{tikzpicture}[x=1cm, y=0.7cm,
    hd/.style={hdr, minimum width=8mm, minimum height=6mm},
    it/.style={cell, minimum width=12mm, minimum height=6mm, font=\small},
    it3/.style={it, fill=currentfill, draw=accentorange, line width=1pt}]
    \foreach \u/\lst in {1/{(2,4),(3,1)}, 2/{(1,4),(3,2)}, 3/{(1,1),(2,2),(4,7)}, 4/{(3,7),(5,3)}, 5/{(4,3)}}{
        \node[hd] (h\u) at (0,-\u) {\u};
        \foreach \e [count=\k] in \lst {
            \ifnum\u=3 \node[it3] (c\u\k) at (0.6+1.35*\k,-\u) {\e};
            \else \node[it] (c\u\k) at (0.6+1.35*\k,-\u) {\e}; \fi
        }
        \draw[flow] (h\u) -- (c\u1);
    }
    \node[font=\scriptsize\bfseries\color{darkgray}] at (0,-0.25) {adj};
    \node[note, anchor=west, text width=5.2cm] at (5.2,-3) {Entry (neighbour, weight). Vertex 3 has degree 3, so its list has 3 entries and ``neighbours of 3'' costs exactly 3 steps.\\[3pt]
        Total length $=\sum\deg(v)=2m=10$ entries. Memory $\bigO{n+m}$.};
\end{tikzpicture}
\figcap{The weighted adjacency list of $G_2$, in input order. Each undirected edge appears twice, once in each endpoint's list (e.g.\ $(3,7)$ in list 4 and $(4,7)$ in list 3).}
\end{diagrambox}

\begin{insightbox}[Why the adjacency list is the default]
A traversal looks at the neighbours of each vertex once. With adjacency lists that costs $\sum_v \deg(v) = 2m$ steps in total (handshake lemma), so the whole traversal is $\bigO{n+m}$. With a matrix the same traversal scans $n$ cells per vertex: $\bigO{n^2}$, even if the graph has only a few edges.
\end{insightbox}

%------------------------------------------------------------
\subsubsection*{4. Forward star and CSR --- adjacency lists in flat arrays}

A \texttt{vector<vector<int>>} allocates one small array per vertex. Two classic layouts store the same lists in a few flat arrays instead --- less memory, fewer allocations, faster on huge inputs ($n,m\approx 10^6$). You will meet both in other people's code.

\textbf{Forward star} (``linked list in arrays'', \texttt{head}/\texttt{nxt}/\texttt{to}). Edges are numbered $0,1,2,\dots$ in the order they are added; \texttt{head[u]} is the last edge added at $u$, and \texttt{nxt[e]} is the edge added at the same vertex before $e$ (or $-1$). Walking \texttt{e = head[u], nxt[e], \ldots} lists $u$'s neighbours in \emph{reverse} input order.

\begin{diagrambox}
\begin{tikzpicture}
    \dpindex{fh}{2.2}{$u$}{1}{5}
    \dparray{hd}{1.5}{head}{2/cell, 4/cell, 6/curcell, 8/cell, 9/cell}
    \begin{scope}[yshift=-0.4cm]
    \dpindex{fe}{0.6}{$e$}{0}{9}
    \dparray{to}{-0.1}{to}{2/cell, 1/cell, 3/cell, 1/cell, 3/cell, 2/cell, 4/curcell, 3/cell, 5/cell, 4/cell}
    \dparray{nx}{-0.8}{nxt}{-1/cell, -1/cell, 0/cell, -1/cell, 1/cell, 3/cell, 5/curcell, -1/cell, 7/cell, -1/cell}
    \end{scope}
    \draw[dep] (hd2.south) to[out=-90, in=90] (to6.north);
    \draw[dep] (nx6.south) to[out=-120, in=-60] (nx5.south);
    \draw[dep] (nx5.south) to[out=-120, in=-60] (nx3.south);
    \node[note, anchor=west, text width=9cm] at (0,-2.75) {Neighbours of 3: \texttt{head[3]}$=6$ gives \texttt{to[6]}$=4$; \texttt{nxt[6]}$=5$ gives \texttt{to[5]}$=2$; \texttt{nxt[5]}$=3$ gives \texttt{to[3]}$=1$; \texttt{nxt[3]}$=-1$: stop. Order: 4, 2, 1.};
\end{tikzpicture}
\figcap{Forward star of $G_2$ (weights omitted; a fourth array \texttt{wt[e]} would hold them). Edge $2i$ is $u_i\to v_i$ and edge $2i+1$ is $v_i\to u_i$, so \texttt{e \^{} 1} is always the reverse edge --- a property flow algorithms rely on.}
\end{diagrambox}

\textbf{CSR} (compressed sparse row). Count the degrees, take prefix sums to get \texttt{start[u]}, then drop each neighbour into its slot. The neighbours of $u$ are \texttt{nbr[start[u]]} \ldots{} \texttt{nbr[start[u+1]-1]}, contiguous in memory and in input order.

\begin{diagrambox}
\begin{tikzpicture}
    \dpindex{ci}{1.4}{$u$}{1}{6}
    \dparray{st}{0.7}{start}{0/cell, 2/cell, 4/curcell, 7/curcell, 9/cell, 10/cell}
    \begin{scope}[yshift=-0.3cm]
    \dpindex{ck}{-0.1}{$k$}{0}{9}
    \dparray{nb}{-0.8}{nbr}{2/cell, 3/cell, 1/cell, 3/cell, 1/curcell, 2/curcell, 4/curcell, 3/cell, 5/cell, 4/cell}
    \end{scope}
    \draw[accentorange, line width=1.3pt, decorate, decoration={brace, amplitude=4pt, mirror}] (3.7,-1.55) -- (6.4,-1.55)
        node[midway, below=4pt, font=\scriptsize\color{accentorange}] {neighbours of 3 = slots $4..6$};
    \draw[dep] (st2.south) to[out=-80, in=100] (nb4.north);
    \draw[dep] (st3.south) to[out=-80, in=100] (nb7.north);
\end{tikzpicture}
\figcap{CSR of $G_2$: \texttt{start} has $n+1$ entries ($u=1..6$ here, with \texttt{start[6]}$=2m=10$ as a sentinel); the slice for vertex 3 runs from \texttt{start[3]}$=4$ to \texttt{start[4]}$-1=6$.}
\end{diagrambox}

\begin{lstlisting}[style=cpp]
#include <bits/stdc++.h>
using namespace std;

// Forward star (head / nxt / to arrays) and CSR (start / nbr arrays)
// built from the same undirected edge list.
const int MAXN = 200005, MAXM = 400010;   // MAXM = 2 * (max edges) + slack
int head[MAXN], nxt[MAXM], to_[MAXM], cnt = 0;

void addEdge(int u, int v) {           // directed u -> v
    to_[cnt] = v;
    nxt[cnt] = head[u];                // new edge points to the old first edge
    head[u] = cnt++;                   // and becomes the first edge of u
}

int main() {
    int n, m;
    if (!(cin >> n >> m)) return 0;
    vector<int> eu(m), ev(m);
    for (int i = 0; i < m; i++) cin >> eu[i] >> ev[i];

    // ---- forward star ----
    fill(head, head + n + 1, -1);
    for (int i = 0; i < m; i++) { addEdge(eu[i], ev[i]); addEdge(ev[i], eu[i]); }
    for (int u = 1; u <= n; u++) {
        cout << "star " << u << ":";
        for (int e = head[u]; e != -1; e = nxt[e]) cout << " " << to_[e];
        cout << "\n";
    }

    // ---- CSR: count, prefix sums, fill ----
    vector<int> start(n + 2, 0);
    for (int i = 0; i < m; i++) { start[eu[i] + 1]++; start[ev[i] + 1]++; }
    for (int u = 1; u <= n + 1; u++) start[u] += start[u - 1];
    vector<int> pos(start), nbr(2 * m);
    for (int i = 0; i < m; i++) { nbr[pos[eu[i]]++] = ev[i]; nbr[pos[ev[i]]++] = eu[i]; }
    for (int u = 1; u <= n; u++) {
        cout << "csr  " << u << ":";
        for (int k = start[u]; k < start[u + 1]; k++) cout << " " << nbr[k];
        cout << "\n";
    }
    return 0;
}
\end{lstlisting}

\begin{lstlisting}[style=trace]
input : 5 5   1 2  1 3  2 3  3 4  4 5          (G2 without weights)
output: star 1: 3 2      csr  1: 2 3
        star 2: 3 1      csr  2: 1 3
        star 3: 4 2 1    csr  3: 1 2 4
        star 4: 5 3      csr  4: 3 5
        star 5: 4        csr  5: 4
\end{lstlisting}
(The program prints the ``star'' block first and the ``csr'' block after it; they are shown side by side here.)

%------------------------------------------------------------
\subsubsection*{Choosing a representation}

\begin{table}[H]
\centering\small
\begin{tabularx}{\linewidth}{@{}X c c c@{}}
\toprule
\textbf{Operation} & \textbf{Edge list} & \textbf{Matrix} & \textbf{Adjacency list / CSR} \\
\midrule
memory & $\bigO{m}$ & $\bigO{n^2}$ & $\bigO{n+m}$ \\
build from input & $\bigO{m}$ & $\bigO{n^2+m}$ & $\bigO{n+m}$ \\
list the neighbours of $u$ & $\bigO{m}$ & $\bigO{n}$ & $\bigO{\deg u}$ \\
is $u$--$v$ an edge? & $\bigO{m}$ & $\bigO{1}$ & $\bigO{\deg u}$ \\
visit every edge from every vertex (a traversal) & --- & $\bigO{n^2}$ & $\bigO{n+m}$ \\
sort edges by weight & $\bigO{m\log m}$ & --- & --- \\
\bottomrule
\end{tabularx}
\end{table}

\begin{warningbox}[{$n = 2\cdot 10^5$: the matrix is impossible}]
An \texttt{int} matrix for $n=2\cdot10^5$ needs $n^2\cdot 4 = 1.6\cdot10^{11}$ bytes $=160$ GB; even one \emph{bit} per cell is 5 GB. A typical memory limit is 256 MB. The adjacency list of the same graph with $m=2\cdot10^5$ edges stores $n+2m=6\cdot10^5$ integers: a few megabytes. Rule of thumb: \textbf{use a matrix only when $n\le 5000$} (an \texttt{int} matrix of $5000^2$ is 100 MB --- already borderline).
\end{warningbox}

%------------------------------------------------------------
\subsubsection*{Reading input correctly}

Four details cause most ``my graph code is right but I get WA/RE'' moments:

\begin{enumerate}[leftmargin=*]
    \item \textbf{0- or 1-indexed?} Codeforces, CSES and AtCoder usually number vertices $1..n$: allocate \texttt{n + 1} and ignore index 0. LeetCode usually numbers them $0..n-1$: allocate \texttt{n}. Mixing them up gives an out-of-bounds access on vertex $n$. Alternatively read \texttt{u}, \texttt{v} and immediately do \texttt{u-{}-, v-{}-}, then work 0-indexed everywhere.
    \item \textbf{Undirected = two directed edges.} Push both directions. Forgetting \texttt{adj[v].push\_back(u)} makes half of the graph invisible; pushing both directions for a directed graph invents edges that do not exist.
    \item \textbf{Multiple test cases.} A global \texttt{adj} must be \emph{cleared} for vertices $1..n$ of every test --- or declared inside the test loop. And clear only what the test uses: resetting a $2\cdot10^5$ array for each of $10^4$ tiny tests is $2\cdot10^9$ operations (\tle). Statements that say ``the sum of $n$ over all tests does not exceed $2\cdot10^5$'' are telling you to allocate per test.
    \item \textbf{Fast I/O.} With $m$ up to $10^6$, \texttt{ios::sync\_with\_stdio(false); cin.tie(nullptr);} is mandatory.
\end{enumerate}

\begin{diagrambox}
\begin{tikzpicture}[x=1cm, y=1cm]
    \node[draw=darkgray, fill=lightgray!60, font=\ttfamily\small, align=left, inner sep=3pt, text width=1.1cm] (t0) at (0,1.6) {2};
    \node[draw=primaryblue, fill=lightgray!60, font=\ttfamily\small, align=left, inner sep=3pt, text width=1.1cm, below=0pt of t0] (ta)
        {5 5\\1 2 4\\1 3 1\\2 3 2\\3 4 7\\4 5 3};
    \node[draw=accentorange, fill=lightgray!60, font=\ttfamily\small, align=left, inner sep=3pt, text width=1.1cm, below=0pt of ta] (tb)
        {3 1\\1 3 5};
    \draw[decorate, decoration={brace, amplitude=4pt}, primaryblue, thick] ([xshift=2pt]ta.north east) -- ([xshift=2pt]ta.south east) node[midway, right=5pt, font=\scriptsize, text=primaryblue] {test 1: $G_2$};
    \draw[decorate, decoration={brace, amplitude=4pt}, accentorange, thick] ([xshift=2pt]tb.north east) -- ([xshift=2pt]tb.south east) node[midway, right=5pt, font=\scriptsize, text=accentorange] {test 2};
    \node[font=\scriptsize\color{darkgray}, left=2pt of t0] {$T$};
    \begin{scope}[xshift=4.6cm, yshift=-0.9cm]
        \node[gv] (1) at (0,0) {1};
        \node[gv] (2) at (1.2,0.9) {2};
        \node[gv] (3) at (1.7,-0.2) {3};
        \draw[ge] (1) -- node[gw, below] {5} (3);
        \node[font=\scriptsize\bfseries, text=accentorange] at (0.85,-1.0) {test 2: $n=3$, $m=1$};
        \node[note, anchor=west, text width=4.6cm] at (2.6,0.4) {Vertex 2 is isolated. If \texttt{adj} were not reset, list 1 would still contain 2 and 3 from test 1 --- a wrong answer that passes test 1 and fails test 2.};
    \end{scope}
\end{tikzpicture}
\figcap{Two tests in one input. Each test must start with empty adjacency lists.}
\end{diagrambox}

Here is a complete reading template that gets all four details right. Copy it, then replace the printing loop with your algorithm.

\begin{lstlisting}[style=cpp]
#include <bits/stdc++.h>
using namespace std;

// Complete reading template: T test cases, each a weighted graph.
// Input:  T
//         n m            (vertices 1..n, m edges)
//         u v w   x m    (undirected edge u-v of weight w)
// Output: for every vertex, its neighbours with weights, and its degree.
int main() {
    ios::sync_with_stdio(false);
    cin.tie(nullptr);
    int T;
    if (!(cin >> T)) return 0;
    while (T--) {
        int n, m;
        cin >> n >> m;
        // fresh containers for every test: nothing leaks between tests
        vector<vector<pair<int,int>>> adj(n + 1);   // adj[u] = {(v, w), ...}
        for (int i = 0; i < m; i++) {
            int u, v, w;
            cin >> u >> v >> w;
            adj[u].push_back(make_pair(v, w));      // u -> v
            adj[v].push_back(make_pair(u, w));      // v -> u (undirected!)
        }
        for (int u = 1; u <= n; u++) {
            cout << u << " (deg " << adj[u].size() << "):";
            for (size_t k = 0; k < adj[u].size(); k++)
                cout << " " << adj[u][k].first << "/" << adj[u][k].second;
            cout << "\n";
        }
        cout << "--\n";
    }
    return 0;
}
\end{lstlisting}

\begin{lstlisting}[style=trace]
output for the input in the figure:
1 (deg 2): 2/4 3/1            1 (deg 1): 3/5
2 (deg 2): 1/4 3/2            2 (deg 0):
3 (deg 3): 1/1 2/2 4/7        3 (deg 1): 1/5
4 (deg 2): 3/7 5/3            --
5 (deg 1): 4/3
--
\end{lstlisting}
(Test 1 on the left, test 2 on the right; the program prints them one after the other.)

\begin{tricksbox}[For a directed graph, and for LeetCode]
For a directed graph delete the second \texttt{push\_back}. On LeetCode the graph arrives as \texttt{int n, vector<vector<int>>\& edges}; the same loop becomes \texttt{for (auto \&e : edges) \{ adj[e[0]].push\_back(e[1]); adj[e[1]].push\_back(e[0]); \}} with \texttt{adj(n)} because vertices are $0..n-1$.
\end{tricksbox}

%------------------------------------------------------------
\subsubsection*{Reading a grid}

Many graphs arrive as a picture: $R$ lines of $C$ characters, where \texttt{.} is free and \texttt{\#} is a wall. We do \emph{not} build an adjacency list for a grid. The cell $(r,c)$ is the vertex, and its neighbours are computed on the fly with two small direction arrays. When an algorithm needs a single integer per vertex (an array index, a DSU element), use $\mathit{id}=r\cdot C+c$.

\begin{diagrambox}
\begin{tikzpicture}[x=0.75cm, y=0.75cm]
    \foreach \row [count=\r from 0] in {{.,.,\#,.},{.,.,.,.},{\#,.,.,\#}}{
        \foreach \ch [count=\c from 0] in \row {
            \pgfmathtruncatemacro{\id}{\r*4+\c}
            \if\ch\#
                \node[gwall] at (\c,-\r) {};
            \else
                \ifnum\id=5 \node[gcell, fill=currentfill] at (\c,-\r) {\scriptsize\id};
                \else \node[gcell] at (\c,-\r) {\scriptsize\id}; \fi
            \fi
        }
    }
    \foreach \c in {0,...,3} \node[font=\scriptsize\color{darkgray}] at (\c,0.75) {$c{=}\c$};
    \foreach \r in {0,...,2} \node[font=\scriptsize\color{darkgray}] at (-0.95,-\r) {$r{=}\r$};
    \foreach \dx/\dy in {0/1, 0/-1, -1/0, 1/0} \draw[dep, line width=1.3pt] (1+0.25*\dx,-1+0.25*\dy) -- (1+0.7*\dx,-1+0.7*\dy);
    \begin{scope}[xshift=4.4cm]
        \node[draw=darkgray, fill=lightgray!60, rounded corners=2pt, font=\ttfamily\small, align=left, inner sep=4pt] at (0.4,-1) {3 4\\..\#.\\....\\\#..\#};
        \node[font=\scriptsize\color{darkgray}] at (0.4,0.45) {input};
    \end{scope}
    \begin{scope}[xshift=7.6cm]
        \foreach \row [count=\r from 0] in {{2,2,\#,1},{2,4,3,2},{\#,2,2,\#}}{
            \foreach \ch [count=\c from 0] in \row {
                \if\ch\# \node[gwall] at (\c,-\r) {};
                \else \node[gcell, font=\small\bfseries] at (\c,-\r) {\ch}; \fi
            }
        }
        \node[font=\scriptsize\color{darkgray}] at (1.5,0.75) {free neighbours of each cell};
    \end{scope}
\end{tikzpicture}
\figcap{Left: cell ids $r\cdot C+c$ for $C=4$; the yellow cell $(1,1)$ has id 5 and four candidate neighbours (up, down, left, right). Right: the output of the program below --- how many of the four candidates are inside the grid and free.}
\end{diagrambox}

\begin{lstlisting}[style=cpp]
#include <bits/stdc++.h>
using namespace std;

// Reading a grid: R rows of C characters, '.' = free, '#' = wall.
// Prints, for every free cell, how many free 4-neighbours it has.
int main() {
    int R, C;
    if (!(cin >> R >> C)) return 0;
    vector<string> g(R);
    for (int r = 0; r < R; r++) cin >> g[r];         // one row = one string
    const int dr[4] = {-1, 1, 0, 0};
    const int dc[4] = {0, 0, -1, 1};
    for (int r = 0; r < R; r++) {
        for (int c = 0; c < C; c++) {
            if (g[r][c] == '#') { cout << '#'; continue; }
            int k = 0;
            for (int d = 0; d < 4; d++) {
                int nr = r + dr[d], nc = c + dc[d];
                if (nr < 0 || nr >= R || nc < 0 || nc >= C) continue;  // off the board
                if (g[nr][nc] == '#') continue;                        // wall
                k++;
            }
            cout << k;
        }
        cout << "\n";
    }
    return 0;
}
\end{lstlisting}

\begin{lstlisting}[style=trace]
input : 3 4  ..#.  ....  #..#
output: 22#1
        2432
        #22#
\end{lstlisting}

The bounds check \emph{before} the wall check matters: \texttt{g[-1][c]} is undefined behaviour. Grids are the subject of Pattern~3.

%------------------------------------------------------------
\subsubsection*{Check your understanding}

\begin{enumerate}[leftmargin=*]
    \item $n=1000$, $m=4000$, and the algorithm asks ``is there an edge $u$--$v$?'' $10^6$ times. Which representation?
    \item $n=10^5$, $m=2\cdot10^5$, undirected. How many entries do the adjacency lists hold in total?
    \item In the forward-star figure, list the neighbours of vertex 2 by following \texttt{head} and \texttt{nxt}.
    \item A LeetCode problem gives \texttt{n = 4, edges = [[0,1],[1,2]]}. What goes wrong with \texttt{vector<vector<int>> adj(n)} if you then write \texttt{adj[u+1]}?
    \item Your solution passes the first test case but fails the second, and the two test cases are each correct when run alone. What is the most likely bug?
\end{enumerate}

\begin{solutionbox}[Answers]
\begin{enumerate}[leftmargin=*]
    \item An adjacency matrix ($10^6$ cells, $\bigO{1}$ per query). Alternatively an adjacency list with each list sorted and binary search, or a hash set of normalised pairs.
    \item $2m = 4\cdot10^5$ entries (each undirected edge is stored twice), plus $n+1$ list headers.
    \item \texttt{head[2]}$=4$: \texttt{to[4]}$=3$; \texttt{nxt[4]}$=1$: \texttt{to[1]}$=1$; \texttt{nxt[1]}$=-1$: stop. Neighbours 3, 1 --- matching ``star 2: 3 1'' in the output.
    \item Vertex 3 becomes index 4, which is out of bounds for a vector of size 4: undefined behaviour (crash or silent corruption). Stay 0-indexed on LeetCode.
    \item State that survives between tests: an adjacency list, \texttt{visited} array or counter that was not reset.
\end{enumerate}
\end{solutionbox}
